% sections/evaluation.tex — Evaluation
\section{Evaluation}
\label{sec:evaluation}

The evaluation is structured around the formal model: each
benchmark either validates a specific operational primitive on
which the model depends, or measures a system-level property the
model predicts.  We do not claim a comprehensive comparison against
all GPU programming abstractions; we benchmark against the
host-driven kernel-launch baseline and against CUDA Graphs (the
strongest publicly-documented host-driven alternative), and
report the headline numbers that the formal model's headline
claims rest on.

All measurements were collected on 2026-04-16 on a single H100~NVL
under the conditions described in
\Cref{sec:eval:methodology}; raw data is available in
\texttt{target/criterion/} in the open-source distribution, and the
methodology document is \texttt{docs/benchmarks/METHODOLOGY.md}.

\subsection{Methodology}
\label{sec:eval:methodology}

We follow the protocol of \citet{moss2002benchmarks}, augmented for
GPU benchmarking with the recommendations of
\citet{georges2007rigorous} on confidence-interval reporting and
between-condition statistical testing.  Concretely:

\begin{itemize}
\item GPU clocks are locked at 1{,}785~MHz via
\texttt{nvidia-smi -lgc}; persistence mode is enabled
(\texttt{nvidia-smi -pm 1}); compute mode is set to
\texttt{EXCLUSIVE\_PROCESS} to preclude background contention.

\item Each measurement is replicated across 20 independent trials,
each consisting of 1{,}000 inner iterations; the per-iteration
timing is reported with a 95\% confidence interval.

\item Between-condition comparisons (e.g., persistent actor vs.\
\texttt{cuLaunchKernel}) report Cohen's~\(d\) and a Welch's
\(t\)-test \(p\)-value; the difference is considered significant if
\(p < 10^{-3}\) and \(|d| > 2\) (very large effect).

\item All measurements are taken on Azure
\texttt{Standard\_NC80adis\_H100\_v5} (AMD EPYC 9V84 80-vCPU,
314~GiB DDR5, 2\(\times\) H100~NVL with 95{,}830~MiB HBM3 each
and ECC enabled, NV12 NVLink topology between GPU0 and GPU1),
CUDA Toolkit
12.8 (V12.8.93), driver 595.58.03.

\item ECC error counts are sampled before and after each trial;
trials in which any ECC corrected error appears are discarded
(zero discards across the reported runs).
\end{itemize}

\subsection{Command Injection}
\label{sec:eval:injection}

The first claim is that mapped-memory command injection eliminates
the kernel-launch overhead that dominates host-orchestrated GPU
workloads.  We compare three execution models (\Cref{tab:injection}):

\begin{table}[h]
\centering
\begin{tabular}{lrrrr}
\toprule
Model & Per-cmd ($\mu$s) & Total ($\mu$s) & 95\% CI ($\mu$s) & vs.\ baseline \\
\midrule
\texttt{cuLaunchKernel} (baseline) & 1.583 & 1{,}583.1 & $\pm$6.0 & 1.0$\times$ \\
CUDA Graph replay & 0.547 & 546.9 & $\pm$12.3 & 2.9$\times$ \\
Persistent actor (mapped mem) & 0.000 & 0.2 & $\pm$0.0 & \textbf{8{,}698$\times$} \\
\bottomrule
\end{tabular}
\caption{Command injection latency, 1{,}000 commands per trial, 20
trials.  Persistent-actor latency is the time to issue a
\texttt{volatile} write to a 32-bit slot in mapped memory; the
device picks up the command at the start of its next tick.}
\label{tab:injection}
\end{table}

The Welch's \(t\)-test on the persistent-actor versus
\texttt{cuLaunchKernel} comparison gives \(p < 10^{-20}\) with
Cohen's \(d \gg 2.0\); the persistent-actor versus CUDA-Graph
comparison is similarly significant (\(p < 10^{-20}\), \(d \gg
2.0\)).  The 8{,}698$\times$ figure is the strongest single result
in the paper: it is what makes the persistent-kernel substrate
viable for sub-microsecond actor injection at all.

The interpretation, in terms of the formal model, is that the
\textsc{Process} step's setup cost (the host's contribution to
``ready the actor for one tick'') is essentially free in the
persistent-kernel model and substantial in either of the
host-orchestrated alternatives.  Workloads in which the actor's
work per tick is small (telemetry aggregation, sensor fusion,
streaming graph updates) can amortize the host's setup cost only in
the persistent-kernel design.

\subsection{K2K Throughput and Stability}
\label{sec:eval:throughput}

The second claim is that the K2K calculus's lock-free queues sustain
high throughput at low latency, with stability sufficient for
operational use.  We measured per-message latency at four payload
sizes (\Cref{tab:k2k-latency}) and a 60-second sustained run at the
1~KiB payload size to characterize stability
(\Cref{tab:k2k-sustained}).

\begin{table}[h]
\centering
\begin{tabular}{lrrr}
\toprule
Payload & Latency (ns) & 95\% CI (ns) & Throughput (M~msg/s) \\
\midrule
64 B   & 72.28  & [72.26, 72.32] & 13.83 \\
256 B  & 74.67  & [74.65, 74.70] & 13.39 \\
1 KiB  & 82.64  & [82.61, 82.68] & 12.10 \\
4 KiB  & 177.50 & [177.47, 177.53] & 5.63 \\
\bottomrule
\end{tabular}
\caption{K2K per-message latency by payload size, 1{,}000 messages
per trial, 20 trials, in-block (\textsc{Smem}) tier.}
\label{tab:k2k-latency}
\end{table}

\begin{table}[h]
\centering
\begin{tabular}{lr}
\toprule
Metric & Value \\
\midrule
Duration per trial & 60.0 s \\
Trials & 4 \\
Mean throughput & \textbf{5.10~M~msg/s} \\
Coefficient of variation (across all 48 windows) & \textbf{0.66\%} \\
p50 latency & 90 ns \\
p95 latency & 100 ns \\
p99 latency & 110 ns \\
Throughput drift (first vs.\ last window per trial) & $<$2\% \\
Memory growth & 0 MiB \\
ECC errors (pre $\to$ post per GPU) & 0 / 0 \\
\bottomrule
\end{tabular}
\caption{K2K sustained run, four 60-second trials at 1~KiB payload
on 2$\times$ H100~NVL.  Per-window samples taken every 5 seconds (12
windows per trial); the CV of 0.66\% across the resulting
48-window pool sits well below the conventional ``extremely
stable'' threshold of 1\%.  Per-window p99 latency is flat at
110~ns across all 48 windows — no thermal throttle, no drift.}
\label{tab:k2k-sustained}
\end{table}

The throughput in \Cref{tab:k2k-sustained} is lower than the
inverse of the 90~ns~p50 latency would suggest: the gap is the
device-side BSP barrier between phases of the persistent kernel
(per \Cref{sec:implementation:kernel}).  Each iteration of the
kernel processes one message per actor; the throughput is the
product of actor count and tick rate.  At 5.10~M~msg/s with a
4-block, 256-thread cluster, this is roughly 5.0~K
messages/actor/second, sustained for 60 seconds without drift.

\paragraph{Per-tier round-trip latency}
The K2K channel hierarchy of \Cref{sec:k2k:tiers} predicts a
clean monotonic ordering \textsc{Smem}~$<$~\textsc{Dsmem}~$<$~\textsc{Hbm}.
\Cref{tab:k2k-tier} measures all three tiers at the same four
payload sizes via a host-observed kernel launch + synchronization
round-trip (the relevant operational cost for the actor system,
which always reaches each tier through a launch).  The per-message
\emph{device-cycle} latency is sub-microsecond on every
tier~\citep{nvidia2022hopper}; what
\Cref{tab:k2k-tier} reports is the host-side wall-clock that the
actor runtime actually pays.

\begin{table}[h]
\centering
\begin{tabular}{lrrrr}
\toprule
Payload & SMEM (ns) & DSMEM (ns) & HBM (ns) & HBM/SMEM \\
\midrule
64 B   & 6{,}793 $\pm$ 41 &  9{,}012 $\pm$ 40 & 10{,}460 $\pm$ 39 & 1.54$\times$ \\
256 B  & 6{,}702 $\pm$ 31 & 10{,}598 $\pm$ 34 & 12{,}066 $\pm$ 33 & 1.80$\times$ \\
1 KiB  & 6{,}722 $\pm$ 34 & 14{,}059 $\pm$ 34 & 15{,}811 $\pm$ 145 & 2.35$\times$ \\
4 KiB  & 6{,}749 $\pm$ 43 & 15{,}289 $\pm$ 106 & 18{,}237 $\pm$ 30 & 2.70$\times$ \\
\bottomrule
\end{tabular}
\caption{K2K tier-by-tier host-observed wall-clock, 1{,}000 trials
per cell, 200 warmup iterations.  SMEM is payload-independent
(state lives in each block's shared memory; the launch cost
dominates).  DSMEM scales with payload as data flows through the
DSMEM addressing path.  HBM adds the global-memory round-trip on
top of DSMEM, growing the HBM/SMEM ratio from 1.54$\times$ at 64~B
to 2.70$\times$ at 4~KiB — clean monotonic hierarchy across every
payload size, validating the calculus in \Cref{sec:k2k:tiers}.}
\label{tab:k2k-tier}
\end{table}

\subsection{Cluster-Level Synchronization}
\label{sec:eval:cluster}

The third claim is that Hopper's cluster-scoped synchronization is
substantially cheaper than grid-wide synchronization, and that
this difference is exploitable by the K2K calculus's tier choice
(\Cref{sec:k2k:tiers}).  \Cref{tab:cluster-sync} compares
\texttt{cluster.sync()} and \texttt{grid.sync()} under identical
conditions.

\begin{table}[h]
\centering
\begin{tabular}{lrrrr}
\toprule
Method & Per-sync ($\mu$s) & Total ($\mu$s) & 95\% CI ($\mu$s) & Std Dev ($\mu$s) \\
\midrule
\texttt{cluster.sync()} & \textbf{0.628} & 628.2 & $\pm$1.4 & 3.1 \\
\texttt{grid.sync()}    & 1.875          & 1{,}874.9 & $\pm$3.5 & 8.0 \\
\bottomrule
\end{tabular}
\caption{Synchronization latency, 1{,}000 syncs per trial, 20
trials, 4-block $\times$ 256-thread workload.}
\label{tab:cluster-sync}
\end{table}

The 2.98$\times$ speedup of \texttt{cluster.sync()} over
\texttt{grid.sync()} (\(p < 10^{-20}\), Cohen's \(d \gg 2.0\))
follows directly from the architecture: cluster-scoped barriers
synchronize only blocks within a single GPC, while grid-scoped
barriers synchronize across all GPCs and incur the GPC-to-GPC
communication cost.  In K2K terms, this means the
\textsc{Smem} and \textsc{Dsmem} tiers can iterate at
0.628~$\mu$s/tick while the \textsc{Hbm} tier is bounded by the
1.875~$\mu$s/tick of \texttt{grid.sync()}.  The K2K tier choice is
not just a memory-tier choice; it is also a synchronization-tier
choice.

\subsection{Asynchronous Memory Allocation}
\label{sec:eval:async-mem}

The fourth claim is that stream-ordered asynchronous allocation
(\Cref{sec:bg:gpu}, third paragraph) eliminates the global driver
lock that dominates the cost of classical \texttt{cuMemAlloc}.
\Cref{tab:async-alloc} reports the comparison at a 256-byte
allocation size.

\begin{table}[h]
\centering
\begin{tabular}{lrr}
\toprule
Method & Per-alloc (ns) & Speedup \\
\midrule
\texttt{cuMemAlloc} (synchronous) & 102{,}576 & 1.0$\times$ \\
\texttt{cuMemAllocAsync} (stream-ordered) & \textbf{878} & \textbf{116.9$\times$} \\
\bottomrule
\end{tabular}
\caption{Allocation latency, 256-byte allocations, 1{,}000
allocations per trial, 20 trials.}
\label{tab:async-alloc}
\end{table}

The 116.9$\times$ improvement is essential to the
``free actor allocation in the actor's own stream'' assumption of
\Cref{sec:implementation:driver}: at 102~$\mu$s per allocation,
spawning a thousand actors would consume 100~ms of host time and
serialize against every other CUDA stream; at 878~ns, the same
thousand spawns consume under 1~ms and overlap freely with other
stream work.

\subsection{Lifecycle Rule Overhead}
\label{sec:eval:lifecycle}

The lifecycle calculus of \Cref{sec:model:lifecycle} prescribes one
SOS rule per supervisor-issued event (\textsc{Spawn}, \textsc{Activate},
\textsc{Quiesce}, \textsc{Terminate}, \textsc{Restart}).  The
host-side cost of each rule is the time from the supervisor writing
the command into mapped memory to reading back the device-issued
acknowledgement.

\begin{table}[h]
\centering
\begin{tabular}{lrrr}
\toprule
Rule & Mean (ns) & Std dev & p99 (ns) \\
\midrule
\textsc{Spawn}     & 23.1 & 4.7 & 30 \\
\textsc{Activate}  & 23.4 & 4.7 & 30 \\
\textsc{Quiesce}   & 23.8 & 4.8 & 30 \\
\textsc{Terminate} & 23.7 & 4.8 & 30 \\
\textsc{Restart}   & 23.1 & 4.6 & 30 \\
\bottomrule
\end{tabular}
\caption{Lifecycle rule wall-clock, 500 trials per rule, 50
warmup iterations.  All five rules cluster around 23~ns mean /
30~ns p99 — flat within measurement noise.  This is the cost
of issuing a single SOS rule from the host; the device-side
processing overlaps with the next supervisor action through the
mapped-memory ring.}
\label{tab:lifecycle}
\end{table}

The 23-ns mean is two orders of magnitude below the
\emph{sub-microsecond} ceiling we set in \Cref{sec:disc:limits}; in
operational terms, a supervisor can issue 40~M~rule events/second
to the device without saturating the mapped-memory channel.  The
cost is the same for every rule because the device-side dispatch
is a single switch on the command opcode — no per-rule state
machinery on the issue path.

\subsection{Multi-GPU NVLink Bandwidth}
\label{sec:eval:mgpu-bw}

For the cross-GPU K2K calculus of \Cref{sec:migration:cross-gpu},
the central question is whether the runtime can saturate NVLink
P2P bandwidth at realistic transfer sizes.  We measured back-to-back
\texttt{cuMemcpyPeerAsync} bursts (256 transfers per measurement,
32 warmup iterations, single stream) on a 2$\times$ H100~NVL
system with the 12-link bonded NVLink bundle, theoretical peak
$\sim$318~GB/s unidirectional.

\begin{table}[h]
\centering
\begin{tabular}{lrrr}
\toprule
Size & Total (ns) & Bandwidth (GB/s) & Peak utilization \\
\midrule
4 KiB   &        455{,}097 &   2.30 &  0.7\% \\
64 KiB  &        535{,}035 &  31.36 &  9.9\% \\
1 MiB   &      1{,}500{,}029 & 178.95 & 56.3\% \\
16 MiB  &     16{,}661{,}095 & 257.78 & \textbf{81.1\%} \\
\bottomrule
\end{tabular}
\caption{Multi-GPU K2K sustained bandwidth via NVLink P2P,
2$\times$ H100~NVL on Azure NC80adis\_H100\_v5 (NV12 topology).
At 16~MiB the runtime reaches 81\% of the NVLink bundle's
theoretical peak — the stream is saturating the fabric.  The
small-payload path is dominated by per-call overhead, not
bandwidth.  This validates the cross-GPU K2K design: the
inter-GPU tier is bandwidth-bound, not setup-bound, at the
sizes where it matters operationally.}
\label{tab:mgpu-bw}
\end{table}

A complementary one-shot latency measurement at the same sizes
(\Cref{tab:migration} below) characterizes the small-payload
regime: NVLink P2P is 1.96$\times$ faster than host-staged at
1~KiB, growing to 8.65$\times$ at 16~MiB.  The two tables
together cover the bandwidth/latency Pareto front for
cross-GPU actor migration.

\begin{table}[h]
\centering
\begin{tabular}{lrrr}
\toprule
Size   & NVLink p99 (ns) & Host p99 (ns) & P2P speedup (mean) \\
\midrule
1 KiB   &      8{,}029 &     13{,}840 & 1.96$\times$ \\
4 KiB   &      7{,}850 &     16{,}770 & 2.02$\times$ \\
64 KiB  &      8{,}180 &     20{,}120 & 2.44$\times$ \\
1 MiB   &     12{,}630 &     57{,}510 & 5.08$\times$ \\
16 MiB  &     71{,}439 &    603{,}605 & \textbf{8.65$\times$} \\
\bottomrule
\end{tabular}
\caption{NVLink P2P vs.\ host-staged migration latency, 200 trials
per cell, 50 warmup iterations.  At small payloads the per-call
driver overhead dominates and NVLink wins by $\sim$2$\times$; at
16~MiB the bandwidth gap dominates and NVLink wins by 8.65$\times$.
The host path's per-MiB cost grows roughly linearly with size
because the data crosses the PCIe + host-memory boundary twice.}
\label{tab:migration}
\end{table}

\Cref{thm:bounded-latency} of \Cref{sec:migration:cross-gpu}
guarantees by construction that NVLink-paired actors never take
the host path; \Cref{tab:migration} quantifies what that
construction is worth, from $\sim$2$\times$ at the K2K scale to
$\sim$9$\times$ at the bulk-state-transfer scale.

\subsection{Application Workloads}
\label{sec:eval:apps}

To validate that the formal model's primitives compose into useful
applications, we measured three end-to-end workloads built on
\system{}.  We do not analyze them in detail here; the numbers are
included to characterize the system at the application level.

\begin{table}[h]
\centering
\begin{tabular}{lrrr}
\toprule
Workload & Backend & Throughput & Speedup vs.\ CPU \\
\midrule
WaveSim3D (FDTD, $64^3$)        & GPU stencil    & 78{,}089 Mcell/s & 217.9$\times$ \\
WaveSim3D (FDTD, $64^3$)        & GPU block actor & 18{,}633 Mcell/s & 52.1$\times$ \\
TxMon (fraud detection)          & CUDA codegen   & 205.31 G~tx/s & 15{,}624$\times$ \\
ProcInt (process mining DFG)     & GPU            & 10.83 M~ev/s & --- \\
\bottomrule
\end{tabular}
\caption{Three application workloads on H100~NVL.  Speedups are
relative to CPU baseline (Rayon, 40 cores).  WaveSim3D is reported
with two backends to show the cost of the actor abstraction over a
hand-tuned stencil kernel: 4.2$\times$ slowdown for the actor model,
which is the price of the lifecycle and supervision machinery.}
\label{tab:apps}
\end{table}

The WaveSim3D row in \Cref{tab:apps} is the most informative number
for this paper: it isolates the cost of the actor abstraction by
comparing the same workload run as a hand-tuned stencil kernel
versus run as a population of block actors.  The 4.2$\times$ slowdown
is the headline ``what does the actor model cost?'' figure.  For
workloads where the abstraction is essential (multi-actor
composition, supervisor-driven restart, cross-tenant isolation),
the cost is paid once for the framework rather than rebuilt per
application.

\subsection{Mechanized Verification Coverage}
\label{sec:eval:tlc}

The four core theorems of \Cref{thm:lifecycle-soundness,%
thm:no-message-loss,thm:hlc-restart-monotone,thm:cross-tenant-isolation}
plus the migration-safety and bounded-latency theorems of
\Cref{sec:migration} are mechanized in six \TLA{} modules under
\texttt{docs/verification/}.  \Cref{tab:tlc-stats} reports the
state-space size and wall-clock for each spec under its bounded
configuration; all six pass with no counterexamples.

\begin{table}[h]
\centering
\begin{tabular}{lrrrrl}
\toprule
Spec & Bounds & Distinct states & Generated & Wall (s) & Result \\
\midrule
\texttt{actor\_lifecycle} & $|A|\!\le\!3,$ MaxSteps$\!\le\!8$              &   2{,}681 &   7{,}015 & $<$1 & OK \\
\texttt{hlc}              & $|N|\!\le\!2,$ MaxPhys$\!\le\!3$               & 4{,}600{,}000 & 16{,}700{,}000 &   29 & OK \\
\texttt{k2k\_delivery}    & $|A|\!\le\!2,$ MaxMsgs$\!\le\!4,$ QCap$\!\le\!2$ &     241 &     399 & $<$1 & OK \\
\texttt{tenant\_isolation}& $|T|\!\le\!2,$ $|K|\!\le\!3,$ MaxMsgs$\!\le\!3$ & 12{,}261 & 26{,}953 & $<$1 & OK \\
\texttt{migration}        & MaxMsgs$\!\le\!4$                              &     725 &   1{,}445 & $<$1 & OK \\
\texttt{multi\_gpu\_k2k}  & $|G|\!\le\!2,$ $|A|\!\le\!3,$ MaxMsgs$\!\le\!3$ &   3{,}452 &   9{,}571 & $<$1 & OK \\
\bottomrule
\end{tabular}
\caption{TLC model-checking statistics for the six \TLA{}
specifications.  Run on TLC 2.19 (Java 17, 8 worker threads on
the 80-vCPU EPYC host).  All specs pass under the bounded
configurations listed; no invariant violations were observed.
The \texttt{hlc} spec has the largest state space because the
Hybrid Logical Clock state vector is structurally larger than the
other modules' state.}
\label{tab:tlc-stats}
\end{table}

\subsection{Reproducibility}
\label{sec:eval:repro}

All measurements are reproducible from the open-source distribution
of \system{} at the commit recorded in
\texttt{docs/benchmarks/h100-b200-baseline.md} (commit 42724ae,
RingKernel 0.4.2; the equivalent v1.0 release artifact is on the
release branch).  The full benchmark suite is invokable as
\texttt{cargo bench --workspace -p ringkernel-cuda --features cuda},
which produces the Criterion HTML reports under
\texttt{target/criterion/}; the methodology document
\texttt{docs/benchmarks/METHODOLOGY.md} describes the GPU clock
locking, exclusive compute mode, and ECC monitoring protocol step
by step, and \texttt{docs/benchmarks/h100-b200-baseline.md} captures
the full system state at the moment of the reported measurements.
